\part{I–II gimnazijos klasė}

\section{Realieji skaičiai ir intervalai}\label{sec:realieji-skaiciai-ir-intervalai}

\subsection{Apibrėžtys}

\begin{apibreztis}
Realiųjų skaičių aibė $\mathbb{R}$ yra racionaliųjų ir iracionaliųjų skaičių aibė. Racionalieji skaičiai užrašomi trupmena $\frac{m}{n}$, kur $m\in\mathbb{Z}$, $n\in\mathbb{N}$, o iracionalieji tokio užrašo neturi.
\end{apibreztis}

\begin{apibreztis}
Intervalas yra realiųjų skaičių aibė, apibrėžta nelygybėmis. Pavyzdžiui, $(a;b)=\{x\in\mathbb{R}:a<x<b\}$, $[a;b]=\{x\in\mathbb{R}:a\le x\le b\}$.
\end{apibreztis}

\begin{apibreztis}
Skaičiaus $x$ modulis yra atstumas nuo $x$ iki nulio skaičių tiesėje:
\[
|x|=\begin{cases}x,&x\ge 0,\\-x,&x<0.\end{cases}
\]
\end{apibreztis}

\subsection{Teoremos ir savybės}

\begin{savybe}
Jei $a<b$, tai prie abiejų nelygybės pusių pridėjus tą patį skaičių gaunama ekvivalenti nelygybė: $a+c<b+c$.
\end{savybe}

\begin{savybe}
Jei $a<b$ ir $c>0$, tai $ac<bc$. Jei $c<0$, tai $ac>bc$.
\end{savybe}

\begin{teorema}
Nelygybė $|x-a|<r$, kai $r>0$, yra ekvivalenti intervalui $a-r<x<a+r$.
\end{teorema}
\begin{proof}
Pagal modulio apibrėžimą $|x-a|$ yra atstumas tarp $x$ ir $a$. Sąlyga $|x-a|<r$ reiškia, kad šis atstumas mažesnis už $r$, todėl $x$ yra tarp $a-r$ ir $a+r$. Atvirkščiai, jei $a-r<x<a+r$, tai atstumas nuo $x$ iki $a$ mažesnis už $r$, vadinasi, $|x-a|<r$.
\end{proof}

\subsection{Taisyklės ir algoritmai}

\begin{enumerate}
\item Sprendžiant tiesinę nelygybę, atliekami tie patys veiksmai kaip lygtyse, bet dauginant arba dalijant iš neigiamo skaičiaus nelygybės ženklas pakeičiamas priešingu.
\item Intervalų sankirta randama paliekant tik tuos skaičius, kurie priklauso visiems intervalams.
\item Intervalų sąjunga randama sujungiant visus skaičius, kurie priklauso bent vienam intervalui.
\end{enumerate}

\subsection{Iliustruojantys pavyzdžiai}

\paragraph{1 pavyzdys.}
Išspręskime nelygybę $-3(2x-1)\le 15$.
\[
-6x+3\le 15,\qquad -6x\le 12.
\]
Dalyba iš neigiamo skaičiaus pakeičia ženklą:
\[
x\ge -2.
\]
Sprendinių aibė yra $[-2;+\infty)$. \textbf{Atsakymas: $x\in[-2;+\infty)$.}

\paragraph{2 pavyzdys.}
Raskime intervalų $[-4;3)$ ir $(1;7]$ sankirtą.
Skaičius turi tenkinti abi sąlygas: $-4\le x<3$ ir $1<x\le 7$. Bendra dalis prasideda po $1$ ir baigiasi prieš $3$, todėl
\[
[-4;3)\cap(1;7]=(1;3).
\]
\textbf{Atsakymas: $(1;3)$.}

\paragraph{3 pavyzdys.}
Išspręskime $|x-2|\le 5$.
Pagal modulio nelygybės savybę
\[
-5\le x-2\le 5.
\]
Pridedame $2$:
\[
-3\le x\le 7.
\]
\textbf{Atsakymas: $x\in[-3;7]$.}

\subsection{Ryšys su kitomis temomis}

Realieji skaičiai remiasi racionaliųjų skaičių ir laipsnių samprata iš \ref{sec:racionalieji-skaiciai-ir-laipsniai} bei šaknų tema \ref{sec:saknys-ir-realiojo-skaiciaus-samprata}. Intervalai būtini kvadratinių nelygybių sprendiniams aprašyti \ref{sec:kvadratines-lygtys-ir-nelygybes}, funkcijų apibrėžimo sritims nusakyti \ref{sec:funkcijos-ir-ju-savybes} ir vėlesnei ribų teorijai \ref{sec:ribos-ir-tolydumas}.

\section{Laipsniai, šaknys ir logaritmų pradžia}\label{sec:laipsniai-saknys-ir-logaritmu-pradzia}

\subsection{Apibrėžtys}

\begin{apibreztis}
Kai $a\in\mathbb{R}$ ir $n\in\mathbb{N}$, laipsnis $a^n$ yra sandauga $a\cdot a\cdot\ldots\cdot a$, kurioje yra $n$ vienodų dauginamųjų.
\end{apibreztis}

\begin{apibreztis}
Skaičiaus $a\ge 0$ aritmetinė kvadratinė šaknis yra toks neneigiamas skaičius $\sqrt a$, kad $(\sqrt a)^2=a$.
\end{apibreztis}

\begin{apibreztis}
Kai $a>0$, $a\ne 1$, $b>0$, logaritmas $\log_a b$ yra toks skaičius $x$, kad $a^x=b$.
\end{apibreztis}

\subsection{Teoremos ir savybės}

\begin{savybe}
Jei veiksmai apibrėžti, tai $a^m\cdot a^n=a^{m+n}$, $\frac{a^m}{a^n}=a^{m-n}$, $(a^m)^n=a^{mn}$.
\end{savybe}

\begin{savybe}
Kai $a\ge 0$ ir $b\ge 0$, galioja $\sqrt{ab}=\sqrt a\sqrt b$.
\end{savybe}

\begin{teorema}
Kai $a>0$, $a\ne 1$, $x>0$, $y>0$, galioja $\log_a(xy)=\log_a x+\log_a y$.
\end{teorema}
\begin{proof}
Tegu $\log_a x=m$ ir $\log_a y=n$. Tada $x=a^m$, $y=a^n$, todėl $xy=a^{m+n}$. Pagal logaritmo apibrėžimą $\log_a(xy)=m+n=\log_a x+\log_a y$.
\end{proof}

\subsection{Taisyklės ir algoritmai}

\begin{enumerate}
\item Laipsnius su tuo pačiu pagrindu dauginame sudėdami rodiklius, o dalijame rodiklius atimdami.
\item Šaknį paprastiname iškeldami iš po šaknies didžiausią kvadratinį daugiklį.
\item Logaritminę lygtį pirmiausia perrašome rodikline forma, visada patikrindami, kad logaritmo argumentas būtų teigiamas.
\end{enumerate}

\subsection{Iliustruojantys pavyzdžiai}

\paragraph{1 pavyzdys.}
Supaprastinkime $\frac{2^5\cdot 2^{-1}}{2^2}$.
\[
\frac{2^5\cdot 2^{-1}}{2^2}=2^{5+(-1)-2}=2^2=4.
\]
\textbf{Atsakymas: $4$.}

\paragraph{2 pavyzdys.}
Supaprastinkime $\sqrt{72}$.
\[
\sqrt{72}=\sqrt{36\cdot 2}=\sqrt{36}\sqrt2=6\sqrt2.
\]
\textbf{Atsakymas: $6\sqrt2$.}

\paragraph{3 pavyzdys.}
Išspręskime $\log_3(x-1)=2$.
Pagal logaritmo apibrėžimą
\[
x-1=3^2=9,
\]
todėl $x=10$. Tikriname sritį: $x-1>0$, kai $x>1$, taigi $10$ tinka. \textbf{Atsakymas: $x=10$.}

\subsection{Ryšys su kitomis temomis}

Ši tema tęsia laipsnius iš \ref{sec:racionalieji-skaiciai-ir-laipsniai} ir šaknų sampratą iš \ref{sec:saknys-ir-realiojo-skaiciaus-samprata}. Ji reikalinga kvadratinėms lygtims \ref{sec:kvadratines-lygtys-ir-nelygybes}, funkcijų nagrinėjimui \ref{sec:funkcijos-ir-ju-savybes} ir vėliau pereina į rodiklines bei logaritmines funkcijas \ref{sec:laipsnines-rodiklines-ir-logaritmines-funkcijos}.

\section{Kvadratinės lygtys ir nelygybės}\label{sec:kvadratines-lygtys-ir-nelygybes}

\subsection{Apibrėžtys}

\begin{apibreztis}
Kvadratinė lygtis yra lygtis $ax^2+bx+c=0$, kur $a\ne 0$.
\end{apibreztis}

\begin{apibreztis}
Kvadratinės lygties diskriminantas yra skaičius $D=b^2-4ac$.
\end{apibreztis}

\begin{apibreztis}
Kvadratinė nelygybė yra nelygybė $ax^2+bx+c>0$, $ax^2+bx+c\ge 0$, $ax^2+bx+c<0$ arba $ax^2+bx+c\le 0$, kur $a\ne 0$.
\end{apibreztis}

\subsection{Teoremos ir savybės}

\begin{teorema}
Kvadratinės lygties $ax^2+bx+c=0$ sprendiniai, kai $D\ge 0$, yra
\[
x_{1,2}=\frac{-b\pm\sqrt D}{2a}.
\]
Kai $D<0$, realiųjų sprendinių nėra.
\end{teorema}
\begin{proof}
Dauginame lygtį iš $4a$ ir papildome iki kvadrato:
\[
4a^2x^2+4abx+4ac=0,\qquad (2ax+b)^2=b^2-4ac=D.
\]
Jei $D\ge 0$, gauname $2ax+b=\pm\sqrt D$, todėl $x=\frac{-b\pm\sqrt D}{2a}$. Jei $D<0$, kvadratas negali būti neigiamas, todėl realiųjų sprendinių nėra.
\end{proof}

\begin{savybe}
Jei kvadratinės funkcijos nuliai yra $x_1$ ir $x_2$, tai reiškinys $ax^2+bx+c$ ženklą keičia šiuose taškuose, kai šaknys skirtingos.
\end{savybe}

\begin{savybe}
Jei $a>0$, parabolės šakos nukreiptos aukštyn; jei $a<0$, žemyn.
\end{savybe}

\subsection{Taisyklės ir algoritmai}

\begin{enumerate}
\item Kvadratinei lygčiai apskaičiuojame $D=b^2-4ac$, tada taikome šaknų formulę arba nustatome, kad sprendinių nėra.
\item Kvadratinę nelygybę sprendžiame radę nulinius taškus ir sudarę ženklų lentelę.
\item Jei kvadratinis reiškinys išskaidomas sandauga, galima naudoti dauginamųjų nulius ir ženklus intervaluose.
\end{enumerate}

\subsection{Iliustruojantys pavyzdžiai}

\paragraph{1 pavyzdys.}
Išspręskime $x^2-5x+6=0$.
\[
D=(-5)^2-4\cdot1\cdot6=25-24=1.
\]
\[
x_{1,2}=\frac{5\pm1}{2},
\]
todėl $x_1=2$, $x_2=3$. \textbf{Atsakymas: $x=2$ arba $x=3$.}

\paragraph{2 pavyzdys.}
Išspręskime $x^2-4x+4<0$.
\[
x^2-4x+4=(x-2)^2.
\]
Kvadratas visada neneigiamas ir lygus nuliui tik kai $x=2$, todėl jis niekada nėra mažesnis už nulį. \textbf{Atsakymas: sprendinių nėra.}

\paragraph{3 pavyzdys.}
Išspręskime $-x^2+3x+4\ge 0$.
Dauginame iš $-1$ ir keičiame ženklą:
\[
x^2-3x-4\le 0,\qquad (x-4)(x+1)\le 0.
\]
Sandauga neigiama arba lygi nuliui tarp šaknų, todėl
\[
-1\le x\le 4.
\]
\textbf{Atsakymas: $x\in[-1;4]$.}

\subsection{Ryšys su kitomis temomis}

Kvadratinės lygtys remiasi algebriniais reiškiniais \ref{sec:algebriniai-reiskiniai-ir-daugianariai}, realiaisiais intervalais \ref{sec:realieji-skaiciai-ir-intervalai} ir tiesinių nelygybių patirtimi \ref{sec:tiesines-lygtys-nelygybes-ir-sistemos}. Jos tiesiogiai siejamos su kvadratine funkcija \ref{sec:kvadratine-funkcija} ir vėlesniu optimizavimu \ref{sec:optimizavimas-ir-modeliavimas}.

\section{Funkcijos ir jų savybės}\label{sec:funkcijos-ir-ju-savybes}

\subsection{Apibrėžtys}

\begin{apibreztis}
Funkcija yra taisyklė, kiekvienam apibrėžimo srities elementui $x$ priskirianti tiksliai vieną reikšmę $y=f(x)$.
\end{apibreztis}

\begin{apibreztis}
Funkcijos apibrėžimo sritis $D(f)$ yra visų leistinų argumentų aibė, o reikšmių sritis $E(f)$ yra visų funkcijos reikšmių aibė.
\end{apibreztis}

\begin{apibreztis}
Funkcija vadinama didėjančia intervale, jei iš $x_1<x_2$ seka $f(x_1)\le f(x_2)$, ir mažėjančia, jei iš $x_1<x_2$ seka $f(x_1)\ge f(x_2)$.
\end{apibreztis}

\subsection{Teoremos ir savybės}

\begin{savybe}
Funkcijos grafiko susikirtimai su $x$ ašimi yra lygties $f(x)=0$ sprendiniai.
\end{savybe}

\begin{savybe}
Funkcijos grafiko susikirtimas su $y$ ašimi, jei jis yra, turi ordinatę $f(0)$.
\end{savybe}

\begin{teorema}
Jei funkcija intervale yra griežtai didėjanti arba griežtai mažėjanti, tai kiekvieną savo reikšmę tame intervale ji įgyja ne daugiau kaip vieną kartą.
\end{teorema}
\begin{proof}
Tegu $x_1<x_2$. Jei funkcija griežtai didėjanti, tai $f(x_1)<f(x_2)$; jei griežtai mažėjanti, tai $f(x_1)>f(x_2)$. Abiem atvejais skirtingiems argumentams negali atitikti ta pati reikšmė.
\end{proof}

\subsection{Taisyklės ir algoritmai}

\begin{enumerate}
\item Apibrėžimo sričiai rasti pašaliname reikšmes, dėl kurių vardiklis lygus nuliui, šaknis iš neigiamo skaičiaus arba logaritmo argumentas neteigiamas.
\item Nuliniams taškams rasti sprendžiame lygtį $f(x)=0$.
\item Funkcijos kitimo pobūdį nustatome iš grafiko arba lygindami $f(x_1)$ ir $f(x_2)$ pasirinktame intervale.
\end{enumerate}

\subsection{Iliustruojantys pavyzdžiai}

\paragraph{1 pavyzdys.}
Raskime funkcijos $f(x)=\frac{2x+1}{x-3}$ apibrėžimo sritį.
Vardiklis negali būti lygus nuliui:
\[
x-3\ne 0,\qquad x\ne 3.
\]
Todėl
\[
D(f)=\mathbb{R}\setminus\{3\}.
\]
\textbf{Atsakymas: $D(f)=\mathbb{R}\setminus\{3\}$.}

\paragraph{2 pavyzdys.}
Raskime funkcijos $f(x)=x^2-9$ nulinius taškus.
\[
x^2-9=0,\qquad (x-3)(x+3)=0.
\]
Gauname $x=3$ arba $x=-3$. \textbf{Atsakymas: nuliniai taškai $-3$ ir $3$.}

\paragraph{3 pavyzdys.}
Nustatykime, ar $f(x)=2x-5$ didėja.
Jei $x_1<x_2$, tai
\[
f(x_2)-f(x_1)=(2x_2-5)-(2x_1-5)=2(x_2-x_1)>0.
\]
Vadinasi, $f(x_2)>f(x_1)$, todėl funkcija griežtai didėja visoje $\mathbb{R}$. \textbf{Atsakymas: funkcija griežtai didėja.}

\subsection{Ryšys su kitomis temomis}

Funkcijos pratęsia tiesinės funkcijos temą \ref{sec:tiesine-funkcija-ir-grafikas} ir naudoja realiųjų skaičių intervalus \ref{sec:realieji-skaiciai-ir-intervalai}. Jos būtinos kvadratinei funkcijai \ref{sec:kvadratine-funkcija}, trigonometrinėms funkcijoms \ref{sec:trigonometrines-funkcijos} ir išvestinės taikymams \ref{sec:isvestines-taikymai}.

\section{Kvadratinė funkcija}\label{sec:kvadratine-funkcija}

\subsection{Apibrėžtys}

\begin{apibreztis}
Kvadratinė funkcija yra funkcija $f(x)=ax^2+bx+c$, kur $a\ne 0$.
\end{apibreztis}

\begin{apibreztis}
Kvadratinės funkcijos grafikas yra parabolė. Jos viršūnės abscisė yra $x_0=-\frac{b}{2a}$, ordinatė $y_0=f(x_0)$.
\end{apibreztis}

\begin{apibreztis}
Parabolės simetrijos ašis yra tiesė $x=x_0$, einanti per viršūnę.
\end{apibreztis}

\subsection{Teoremos ir savybės}

\begin{savybe}
Jei $a>0$, parabolė turi mažiausią reikšmę viršūnėje; jei $a<0$, didžiausią reikšmę viršūnėje.
\end{savybe}

\begin{savybe}
Kvadratinės funkcijos nuliai yra kvadratinės lygties $ax^2+bx+c=0$ sprendiniai.
\end{savybe}

\begin{teorema}
Kvadratinė funkcija $f(x)=ax^2+bx+c$ gali būti užrašyta viršūnės forma
\[
f(x)=a(x-x_0)^2+y_0,
\]
kur $x_0=-\frac{b}{2a}$ ir $y_0=f(x_0)$.
\end{teorema}
\begin{proof}
Papildome iki kvadrato:
\[
ax^2+bx+c=a\left(x^2+\frac ba x\right)+c
=a\left(x+\frac{b}{2a}\right)^2+c-\frac{b^2}{4a}.
\]
Kadangi $x_0=-\frac{b}{2a}$, gautas užrašas yra $a(x-x_0)^2+y_0$, kur $y_0=f(x_0)$.
\end{proof}

\subsection{Taisyklės ir algoritmai}

\begin{enumerate}
\item Viršūnei rasti skaičiuojame $x_0=-\frac{b}{2a}$ ir $y_0=f(x_0)$.
\item Grafiko kryptį nustatome pagal koeficientą $a$.
\item Reikšmių sritį nustatome pagal viršūnės reikšmę ir šakų kryptį.
\end{enumerate}

\subsection{Iliustruojantys pavyzdžiai}

\paragraph{1 pavyzdys.}
Raskime $f(x)=x^2-6x+8$ viršūnę.
\[
x_0=-\frac{-6}{2\cdot1}=3,\qquad y_0=f(3)=9-18+8=-1.
\]
\textbf{Atsakymas: viršūnė $(3;-1)$.}

\paragraph{2 pavyzdys.}
Nustatykime $f(x)=-2x^2+4x+1$ didžiausią reikšmę.
\[
x_0=-\frac{4}{2\cdot(-2)}=1.
\]
\[
f(1)=-2+4+1=3.
\]
Kadangi $a<0$, ši reikšmė yra didžiausia. \textbf{Atsakymas: didžiausia reikšmė $3$.}

\paragraph{3 pavyzdys.}
Užrašykime $f(x)=x^2+4x+7$ viršūnės forma.
\[
x^2+4x+7=(x^2+4x+4)+3=(x+2)^2+3.
\]
\textbf{Atsakymas: $f(x)=(x+2)^2+3$.}

\subsection{Ryšys su kitomis temomis}

Kvadratinė funkcija jungia funkcijų savybes \ref{sec:funkcijos-ir-ju-savybes} ir kvadratines lygtis \ref{sec:kvadratines-lygtys-ir-nelygybes}. Parabolės ekstremumai parengia optimizavimo idėjas \ref{sec:optimizavimas-ir-modeliavimas}, o grafiko analizė siejasi su išvestinės taikymais \ref{sec:isvestines-taikymai}.

\section{Lygčių ir nelygybių sistemos}\label{sec:lygciu-ir-nelygybiu-sistemos}

\subsection{Apibrėžtys}

\begin{apibreztis}
Lygčių sistema yra kelios lygtys, kurias tie patys nežinomųjų dydžiai turi tenkinti vienu metu.
\end{apibreztis}

\begin{apibreztis}
Nelygybių sistema yra kelios nelygybės, kurių bendrą sprendinių aibę sudaro visų nelygybių sprendinių sankirta.
\end{apibreztis}

\begin{apibreztis}
Sistemos sprendinys su dviem nežinomaisiais yra sutvarkytoji pora $(x;y)$, tenkinanti kiekvieną sistemos lygtį arba nelygybę.
\end{apibreztis}

\subsection{Teoremos ir savybės}

\begin{savybe}
Dviejų tiesinių lygčių su dviem nežinomaisiais sistema gali turėti vieną sprendinį, be galo daug sprendinių arba neturėti sprendinių.
\end{savybe}

\begin{savybe}
Grafiškai dviejų lygčių sistemos sprendiniai yra atitinkamų grafikų bendri taškai.
\end{savybe}

\begin{teorema}
Jei tiesinių lygčių sistemoje vieną lygtį pakeičiame jos suma su kitos lygties kartotiniu, sprendinių aibė nepasikeičia.
\end{teorema}
\begin{proof}
Jei pora tenkina abi pradines lygtis, ji tenkina ir jų tiesinį derinį. Atvirkščiai, iš naujos lygties atėmus tą patį kitos lygties kartotinį atkuriama pradinė lygtis. Todėl sprendinių aibės sutampa.
\end{proof}

\subsection{Taisyklės ir algoritmai}

\begin{enumerate}
\item Keitimo būdu iš vienos lygties išreiškiame vieną nežinomąjį ir įstatome į kitą lygtį.
\item Sudėties būdu lygtys dauginamos taip, kad sudėjus išnyktų vienas nežinomasis.
\item Nelygybių sistemoje kiekvieną nelygybę sprendžiame atskirai, o atsakymui imame sprendinių sankirtą.
\end{enumerate}

\subsection{Iliustruojantys pavyzdžiai}

\paragraph{1 pavyzdys.}
Išspręskime sistemą
\[
\begin{cases}
x+y=7,\\
x-y=1.
\end{cases}
\]
Sudedame lygtis:
\[
2x=8,\qquad x=4.
\]
Tada $4+y=7$, todėl $y=3$. \textbf{Atsakymas: $(4;3)$.}

\paragraph{2 pavyzdys.}
Išspręskime sistemą
\[
\begin{cases}
y=2x+1,\\
y=x^2-1.
\end{cases}
\]
Sulyginame:
\[
2x+1=x^2-1,\qquad x^2-2x-2=0.
\]
\[
x=\frac{2\pm\sqrt{4+8}}2=1\pm\sqrt3.
\]
Tada $y=2x+1$, todėl $y=3\pm2\sqrt3$. \textbf{Atsakymas: $(1+\sqrt3;3+2\sqrt3)$ ir $(1-\sqrt3;3-2\sqrt3)$.}

\paragraph{3 pavyzdys.}
Išspręskime nelygybių sistemą
\[
\begin{cases}
x-2>0,\\
3x+1\le 10.
\end{cases}
\]
Pirmoji duoda $x>2$, antroji $3x\le 9$, taigi $x\le 3$. Sankirta:
\[
2<x\le 3.
\]
\textbf{Atsakymas: $x\in(2;3]$.}

\subsection{Ryšys su kitomis temomis}

Sistemos remiasi tiesinėmis lygtimis ir sistemomis iš \ref{sec:tiesines-lygtys-nelygybes-ir-sistemos}, intervalais \ref{sec:realieji-skaiciai-ir-intervalai} ir funkcijų grafikais \ref{sec:funkcijos-ir-ju-savybes}. Toliau jos plečiamos matricų metodu \ref{sec:matricos-ir-tiesines-sistemos}.

\section{Trigonometrija stačiajame trikampyje}\label{sec:trigonometrija-staciajame-trikampyje}

\subsection{Apibrėžtys}

\begin{apibreztis}
Stačiojo trikampio smailiojo kampo sinusas yra priešingojo statinio ir įžambinės santykis.
\end{apibreztis}

\begin{apibreztis}
Stačiojo trikampio smailiojo kampo kosinusas yra gretimojo statinio ir įžambinės santykis.
\end{apibreztis}

\begin{apibreztis}
Stačiojo trikampio smailiojo kampo tangentas yra priešingojo ir gretimojo statinių santykis.
\end{apibreztis}

\subsection{Teoremos ir savybės}

\begin{savybe}
Kiekvienam smailiajam kampui $\alpha$ galioja $0<\sin\alpha<1$ ir $0<\cos\alpha<1$.
\end{savybe}

\begin{teorema}
Kiekvienam smailiajam kampui $\alpha$ galioja tapatybė $\sin^2\alpha+\cos^2\alpha=1$.
\end{teorema}
\begin{proof}
Tegu stačiojo trikampio statiniai yra $a$ ir $b$, įžambinė $c$, o kampui $\alpha$ priešingas statinis yra $a$. Tada $\sin\alpha=\frac ac$, $\cos\alpha=\frac bc$. Pagal Pitagoro teoremą $a^2+b^2=c^2$, todėl
\[
\sin^2\alpha+\cos^2\alpha=\frac{a^2}{c^2}+\frac{b^2}{c^2}=\frac{c^2}{c^2}=1.
\]
\end{proof}

\begin{savybe}
Kai $\alpha$ smailusis kampas, $\tan\alpha=\frac{\sin\alpha}{\cos\alpha}$.
\end{savybe}

\subsection{Taisyklės ir algoritmai}

\begin{enumerate}
\item Pirmiausia pažymime kampą ir nustatome priešingąjį statinį, gretimąjį statinį bei įžambinę.
\item Pasirenkame santykį $\sin$, $\cos$ arba $\tan$ pagal tai, kurie dydžiai žinomi ir kurį reikia rasti.
\item Jei reikia kampo dydžio, naudojame atvirkštines trigonometrines funkcijas arba lenteles.
\end{enumerate}

\subsection{Iliustruojantys pavyzdžiai}

\paragraph{1 pavyzdys.}
Stačiajame trikampyje kampui $\alpha$ priešingas statinis lygus $6$, o įžambinė $10$. Raskime $\sin\alpha$.
\[
\sin\alpha=\frac{\text{priešingasis statinis}}{\text{įžambinė}}=\frac{6}{10}=\frac35.
\]
\textbf{Atsakymas: $\sin\alpha=\frac35$.}

\paragraph{2 pavyzdys.}
Stačiajame trikampyje $\cos\alpha=\frac45$, o įžambinė lygi $15$. Raskime gretimąjį statinį.
\[
\frac{\text{gretimasis statinis}}{15}=\frac45,
\]
todėl gretimasis statinis lygus
\[
15\cdot\frac45=12.
\]
\textbf{Atsakymas: $12$.}

\paragraph{3 pavyzdys.}
Jei $\sin\alpha=\frac{5}{13}$, raskime $\cos\alpha$, kai $\alpha$ smailusis.
\[
\sin^2\alpha+\cos^2\alpha=1,\qquad \cos^2\alpha=1-\frac{25}{169}=\frac{144}{169}.
\]
Kadangi $\alpha$ smailusis, $\cos\alpha>0$, todėl
\[
\cos\alpha=\frac{12}{13}.
\]
\textbf{Atsakymas: $\cos\alpha=\frac{12}{13}$.}

\subsection{Ryšys su kitomis temomis}

Trigonometrija remiasi Pitagoro teorema \ref{sec:pitos-teorema}, panašiaisiais trikampiais \ref{sec:panasus-trikampiai-ir-talio-teorema} ir funkcijų idėja \ref{sec:funkcijos-ir-ju-savybes}. Vėliau ji pereina į trigonometrines funkcijas \ref{sec:trigonometrines-funkcijos} ir trigonometrines lygtis \ref{sec:trigonometrines-lygtys-ir-tapatybes}.

\section{Apskritimo geometrija}\label{sec:apskritimo-geometrija}

\subsection{Apibrėžtys}

\begin{apibreztis}
Apskritimas yra visų plokštumos taškų, nutolusių nuo centro tuo pačiu atstumu, aibė. Tas atstumas vadinamas spinduliu.
\end{apibreztis}

\begin{apibreztis}
Styga yra atkarpa, jungianti du apskritimo taškus. Skersmuo yra styga, einanti per apskritimo centrą.
\end{apibreztis}

\begin{apibreztis}
Liestinė yra tiesė, turinti su apskritimu lygiai vieną bendrą tašką.
\end{apibreztis}

\subsection{Teoremos ir savybės}

\begin{teorema}
Apskritimo liestinė lietimosi taške yra statmena spinduliui, nubrėžtam į tą tašką.
\end{teorema}
\begin{proof}
Tegu tiesė liečia apskritimą taške $T$, o $O$ yra centras. Jei atstumas nuo $O$ iki liestinės būtų mažesnis už $OT$, tiesė kirstų apskritimą dviejuose taškuose. Kadangi ji turi tik vieną bendrą tašką, trumpiausia atkarpa nuo $O$ iki tiesės yra $OT$. Trumpiausia atkarpa iki tiesės yra statmuo, todėl $OT$ statmena liestinei.
\end{proof}

\begin{savybe}
Įbrėžtinis kampas, remiantis tuo pačiu lanku kaip centrinis kampas, yra lygus pusei centrinio kampo.
\end{savybe}

\begin{savybe}
Vienodo ilgio stygos yra vienodai nutolusios nuo apskritimo centro.
\end{savybe}

\subsection{Taisyklės ir algoritmai}

\begin{enumerate}
\item Uždaviniuose su liestine visada verta nubrėžti spindulį į lietimosi tašką ir pažymėti statų kampą.
\item Įbrėžtinius kampus lyginame pagal lankus, į kuriuos jie remiasi.
\item Apskritimo ilgiui ir skritulio plotui taikome $C=2\pi r$ ir $S=\pi r^2$.
\end{enumerate}

\subsection{Iliustruojantys pavyzdžiai}

\paragraph{1 pavyzdys.}
Apskritimo spindulys $r=7$. Raskime apskritimo ilgį.
\[
C=2\pi r=2\pi\cdot7=14\pi.
\]
\textbf{Atsakymas: $14\pi$.}

\paragraph{2 pavyzdys.}
Centrinis kampas, remiantis lanku $AB$, lygus $80^\circ$. Raskime įbrėžtinį kampą, kuris remiasi tuo pačiu lanku.
\[
\angle ACB=\frac{80^\circ}{2}=40^\circ.
\]
\textbf{Atsakymas: $40^\circ$.}

\paragraph{3 pavyzdys.}
Iš taško $P$ nubrėžta liestinė $PT$ į apskritimą. $OP=13$, $OT=5$. Raskime $PT$.
Kadangi $OT\perp PT$, trikampis $OPT$ statusis:
\[
PT^2=OP^2-OT^2=13^2-5^2=169-25=144.
\]
\[
PT=12.
\]
\textbf{Atsakymas: $12$.}

\subsection{Ryšys su kitomis temomis}

Apskritimo geometrija pratęsia apskritimo ir skritulio temą \ref{sec:apskritimas-skritulys-ir-pi}, remiasi Pitagoro teorema \ref{sec:pitos-teorema} ir kampų savybėmis \ref{sec:kampai-trikampiai-ir-keturkampiai}. Vėliau apskritimai siejami su analitine geometrija \ref{sec:analitine-geometrija-ir-vektoriai}.

\section{Vektoriai ir koordinatės}\label{sec:vektoriai-ir-koordinates}

\subsection{Apibrėžtys}

\begin{apibreztis}
Vektorius yra kryptinė atkarpa, turinti ilgį ir kryptį.
\end{apibreztis}

\begin{apibreztis}
Jei $A(x_1;y_1)$ ir $B(x_2;y_2)$, tai vektoriaus $\overrightarrow{AB}$ koordinatės yra $(x_2-x_1;\,y_2-y_1)$.
\end{apibreztis}

\begin{apibreztis}
Vektoriaus $\vec a=(a_1;a_2)$ ilgis yra $|\vec a|=\sqrt{a_1^2+a_2^2}$.
\end{apibreztis}

\subsection{Teoremos ir savybės}

\begin{savybe}
Vektoriai sudedami sudedant atitinkamas koordinates: $(a_1;a_2)+(b_1;b_2)=(a_1+b_1;\,a_2+b_2)$.
\end{savybe}

\begin{savybe}
Vektorių dauginant iš skaičiaus $k$, kiekviena koordinatė dauginama iš $k$: $k(a_1;a_2)=(ka_1;ka_2)$.
\end{savybe}

\begin{teorema}
Atstumas tarp taškų $A(x_1;y_1)$ ir $B(x_2;y_2)$ yra
\[
AB=\sqrt{(x_2-x_1)^2+(y_2-y_1)^2}.
\]
\end{teorema}
\begin{proof}
Atkarpa $AB$ su horizontaliu ir vertikaliu poslinkiais sudaro statųjį trikampį, kurio statiniai yra $|x_2-x_1|$ ir $|y_2-y_1|$. Pagal Pitagoro teoremą įžambinės ilgis yra nurodyta šaknis.
\end{proof}

\subsection{Taisyklės ir algoritmai}

\begin{enumerate}
\item Norėdami rasti $\overrightarrow{AB}$, iš taško $B$ koordinačių atimame taško $A$ koordinates.
\item Vektorių sumai sudedame koordinates po vieną.
\item Atstumui tarp taškų rasti taikome koordinatinę Pitagoro teoremos formą.
\end{enumerate}

\subsection{Iliustruojantys pavyzdžiai}

\paragraph{1 pavyzdys.}
Raskime $\overrightarrow{AB}$, kai $A(-2;5)$, $B(4;1)$.
\[
\overrightarrow{AB}=(4-(-2);\,1-5)=(6;-4).
\]
\textbf{Atsakymas: $(6;-4)$.}

\paragraph{2 pavyzdys.}
Raskime vektoriaus $\vec a=(3;-4)$ ilgį.
\[
|\vec a|=\sqrt{3^2+(-4)^2}=\sqrt{9+16}=5.
\]
\textbf{Atsakymas: $5$.}

\paragraph{3 pavyzdys.}
Raskime atstumą tarp taškų $P(1;2)$ ir $Q(7;10)$.
\[
PQ=\sqrt{(7-1)^2+(10-2)^2}=\sqrt{36+64}=\sqrt{100}=10.
\]
\textbf{Atsakymas: $10$.}

\subsection{Ryšys su kitomis temomis}

Vektoriai remiasi koordinatėmis iš \ref{sec:dydziu-priklausomybes-ir-koordinaciu-plokstuma} ir Pitagoro teorema \ref{sec:pitos-teorema}. Jie reikalingi funkcijų grafikams \ref{sec:funkcijos-ir-ju-savybes}, analitinei geometrijai \ref{sec:analitine-geometrija-ir-vektoriai} ir erdvės geometrijai \ref{sec:erdves-geometrija}.

\section{Plotai, tūriai ir panašumas}\label{sec:plotai-turiai-ir-panasumas}

\subsection{Apibrėžtys}

\begin{apibreztis}
Figūros plotas yra skaičius, nusakantis, kiek plokštumos dalies užima figūra.
\end{apibreztis}

\begin{apibreztis}
Kūno tūris yra skaičius, nusakantis, kiek erdvės užima kūnas.
\end{apibreztis}

\begin{apibreztis}
Panašios figūros yra vienodos formos figūros, kurių atitinkami kampai lygūs, o atitinkamos kraštinės proporcingos.
\end{apibreztis}

\subsection{Teoremos ir savybės}

\begin{savybe}
Jei panašumo koeficientas yra $k$, tai atitinkamų ilgių santykis lygus $k$.
\end{savybe}

\begin{teorema}
Panašių plokštumos figūrų plotų santykis lygus panašumo koeficiento kvadratui: $\frac{S_2}{S_1}=k^2$.
\end{teorema}
\begin{proof}
Panašioje figūroje kiekvienas ilgis padauginamas iš $k$. Plotą galima matuoti kvadratiniais vienetais, todėl keičiantis abiem kryptims plotas padauginamas iš $k\cdot k=k^2$.
\end{proof}

\begin{savybe}
Panašių erdvinių kūnų tūrių santykis lygus panašumo koeficiento kubui.
\end{savybe}

\subsection{Taisyklės ir algoritmai}

\begin{enumerate}
\item Panašumo koeficientą randame iš atitinkamų kraštinių santykio.
\item Plotų santykiui naudojame $k^2$, o tūrių santykiui $k^3$.
\item Sudėtines figūras skaidome į paprastesnes dalis, kurių plotų ar tūrių formules žinome.
\end{enumerate}

\subsection{Iliustruojantys pavyzdžiai}

\paragraph{1 pavyzdys.}
Du panašūs trikampiai turi atitinkamas kraštines $6$ ir $9$. Mažesniojo plotas $20$. Raskime didesniojo plotą.
\[
k=\frac96=\frac32,\qquad k^2=\frac94.
\]
\[
S_2=20\cdot\frac94=45.
\]
\textbf{Atsakymas: $45$.}

\paragraph{2 pavyzdys.}
Stačiakampio gretasienio matmenys yra $3$, $4$ ir $5$. Raskime tūrį.
\[
V=abc=3\cdot4\cdot5=60.
\]
\textbf{Atsakymas: $60$.}

\paragraph{3 pavyzdys.}
Cilindro spindulys $r=2$, aukštis $h=7$. Raskime tūrį.
\[
V=\pi r^2h=\pi\cdot2^2\cdot7=28\pi.
\]
\textbf{Atsakymas: $28\pi$.}

\subsection{Ryšys su kitomis temomis}

Ši tema remiasi plotais, paviršiais ir tūriais \ref{sec:plotas-pavirsius-ir-turis}, panašiaisiais trikampiais \ref{sec:panasus-trikampiai-ir-talio-teorema} ir erdviniais kūnais \ref{sec:erdviniai-kunai-ir-ju-pjuviai}. Vėliau ji naudojama stereometrijoje \ref{sec:stereometrijos-uzdaviniai} ir integralo taikymuose \ref{sec:integralo-taikymai}.

\section{Statistiniai rodikliai ir koreliacija}\label{sec:statistiniai-rodikliai-ir-koreliacija}

\subsection{Apibrėžtys}

\begin{apibreztis}
Imties vidurkis yra visų duomenų reikšmių suma, padalyta iš jų skaičiaus.
\end{apibreztis}

\begin{apibreztis}
Mediana yra surikiuotos duomenų eilutės vidurinė reikšmė, o kai reikšmių skaičius lyginis, dviejų vidurinių reikšmių vidurkis.
\end{apibreztis}

\begin{apibreztis}
Koreliacija apibūdina dviejų kiekybinių požymių tarpusavio ryšio kryptį ir stiprumą.
\end{apibreztis}

\subsection{Teoremos ir savybės}

\begin{savybe}
Vidurkis jautrus labai didelėms arba labai mažoms reikšmėms, o mediana joms mažiau jautri.
\end{savybe}

\begin{savybe}
Teigiama koreliacija reiškia, kad vienam dydžiui didėjant kitas linkęs didėti; neigiama koreliacija reiškia, kad vienam didėjant kitas linkęs mažėti.
\end{savybe}

\begin{teorema}
Jei prie kiekvienos imties reikšmės pridedame tą patį skaičių $c$, imties vidurkis taip pat padidėja $c$.
\end{teorema}
\begin{proof}
Tegu pradinis vidurkis yra $\bar x=\frac{x_1+\cdots+x_n}{n}$. Naujas vidurkis:
\[
\frac{(x_1+c)+\cdots+(x_n+c)}{n}
=\frac{x_1+\cdots+x_n+nc}{n}
=\bar x+c.
\]
\end{proof}

\subsection{Taisyklės ir algoritmai}

\begin{enumerate}
\item Vidurkiui rasti sudedame visas reikšmes ir dalijame iš jų skaičiaus.
\item Medianai rasti duomenis pirmiausia surikiuojame didėjimo tvarka.
\item Koreliaciją pirmiausia vertiname iš sklaidos diagramos: taškų debesies kryptis ir glaudumas rodo ryšio pobūdį.
\end{enumerate}

\subsection{Iliustruojantys pavyzdžiai}

\paragraph{1 pavyzdys.}
Raskime duomenų $4,6,7,9,14$ vidurkį.
\[
\bar x=\frac{4+6+7+9+14}{5}=\frac{40}{5}=8.
\]
\textbf{Atsakymas: $8$.}

\paragraph{2 pavyzdys.}
Raskime duomenų $12,5,9,7,7,20$ medianą.
Surikiuojame:
\[
5,7,7,9,12,20.
\]
Yra šešios reikšmės, todėl mediana yra dviejų vidurinių vidurkis:
\[
\frac{7+9}{2}=8.
\]
\textbf{Atsakymas: $8$.}

\paragraph{3 pavyzdys.}
Mokinio mokymosi valandos ir testo taškai yra poros $(1;45)$, $(2;55)$, $(3;62)$, $(4;71)$. Nustatykime koreliacijos kryptį.
Didėjant pirmajam dydžiui, antrasis taip pat didėja: $45<55<62<71$. Todėl ryšys yra teigiamas. \textbf{Atsakymas: teigiama koreliacija.}

\subsection{Ryšys su kitomis temomis}

Statistiniai rodikliai tęsia imčių ir sklaidos temą \ref{sec:statistika-imtys-ir-sklaida}. Koreliacija siejasi su funkcijų ir grafikų interpretavimu \ref{sec:funkcijos-ir-ju-savybes}, o vėliau pereina į aprašomąją statistiką ir normalųjį modelį \ref{sec:aprasomoji-statistika-ir-normalusis-modelis} bei statistinę išvadą \ref{sec:statistine-isvada-ir-hipotezes}.

\section{Tikimybė ir kombinatorika}\label{sec:tikimybe-ir-kombinatorika}

\subsection{Apibrėžtys}

\begin{apibreztis}
Baigto bandymo baigčių erdvė $\Omega$ yra visų galimų elementariųjų baigčių aibė.
\end{apibreztis}

\begin{apibreztis}
Įvykis yra baigčių erdvės poaibis. Klasikinė tikimybė lygi $P(A)=\frac{|A|}{|\Omega|}$, kai visos baigtys vienodai tikėtinos.
\end{apibreztis}

\begin{apibreztis}
Kombinatorikoje sandaugos taisyklė teigia: jei pirmą veiksmą galima atlikti $m$ būdų, o po jo antrą $n$ būdų, tai abu veiksmus galima atlikti $mn$ būdų.
\end{apibreztis}

\subsection{Teoremos ir savybės}

\begin{savybe}
Priešingo įvykio tikimybė yra $P(\overline A)=1-P(A)$.
\end{savybe}

\begin{savybe}
Jei įvykiai $A$ ir $B$ nesutaikomi, tai $P(A\cup B)=P(A)+P(B)$.
\end{savybe}

\begin{teorema}
Jei baigtinio bandymo visos baigtys vienodai tikėtinos, tai įvykio ir jo priešingo įvykio tikimybių suma lygi $1$.
\end{teorema}
\begin{proof}
Įvykis $A$ ir priešingas įvykis $\overline A$ neturi bendrų baigčių, o kartu sudaro visą $\Omega$. Todėl $|A|+|\overline A|=|\Omega|$. Padaliję iš $|\Omega|$, gauname $P(A)+P(\overline A)=1$.
\end{proof}

\subsection{Taisyklės ir algoritmai}

\begin{enumerate}
\item Klasikinei tikimybei rasti suskaičiuojame visas vienodai tikėtinas baigtis ir palankias baigtis.
\item Jei reikia bent vieno įvykio, dažnai patogiau skaičiuoti priešingą įvykį.
\item Kombinatoriniuose uždaviniuose atskiriame, ar svarbi tvarka, ar galima kartoti elementus.
\end{enumerate}

\subsection{Iliustruojantys pavyzdžiai}

\paragraph{1 pavyzdys.}
Metamas sąžiningas kauliukas. Raskime tikimybę išmesti lyginį skaičių.
\[
\Omega=\{1,2,3,4,5,6\},\qquad A=\{2,4,6\}.
\]
\[
P(A)=\frac{3}{6}=\frac12.
\]
\textbf{Atsakymas: $\frac12$.}

\paragraph{2 pavyzdys.}
Kiek skirtingų dviženklių skaičių galima sudaryti iš skaitmenų $1,2,3,4$, jei skaitmenys nesikartoja?
Dešimčių skaitmenį galima pasirinkti $4$ būdais, vienetų skaitmenį po to $3$ būdais:
\[
4\cdot3=12.
\]
\textbf{Atsakymas: $12$.}

\paragraph{3 pavyzdys.}
Dėžėje yra $5$ raudoni ir $3$ mėlyni rutuliai. Atsitiktinai traukiamas vienas rutulys. Raskime tikimybę ištraukti ne mėlyną rutulį.
Iš viso yra $8$ rutuliai. Ne mėlyni yra raudoni, jų $5$, todėl
\[
P=\frac58.
\]
\textbf{Atsakymas: $\frac58$.}

\subsection{Ryšys su kitomis temomis}

Tikimybė ir kombinatorika tęsia kombinatorikos pradžią \ref{sec:kombinatorikos-pradzia-ir-tikimybe} ir atsitiktinio bandymo sampratą \ref{sec:atsitiktiniai-bandymas-ir-tikimybe}. Vėliau ši tema pereina į sąlyginę tikimybę \ref{sec:salygine-tikimybe-ir-nepriklausomumas} ir tikimybinius skirstinius \ref{sec:tikimybiniai-skirstiniai}.


\section{Dėsningumai ir sekų apibūdinimas}\label{sec:desningumai-ir-seku-apibudinimas}

\subsection{Apibrėžtys}
\begin{apibreztis}
Tema apima skaičių sekas, nario formulę, rekursinį aprašą ir dėsningumo pagrindimą. Ji nagrinėjama kaip atskira mokymo(si) turinio dalis pagal \emph{matematikaPrograma.pdf}.
\end{apibreztis}
\begin{apibreztis}
Matematinis modelis šioje temoje yra skaičių, simbolių, brėžinių, lentelių arba diagramų sistema, kuria aprašoma reali arba teorinė situacija.
\end{apibreztis}

\subsection{Teoremos ir savybės}
\begin{savybe}
Jeigu uždavinio sąlyga tiksliai perkeliama į matematinį modelį, tai modelio rezultatas gali būti interpretuojamas pradinėje situacijoje.
\end{savybe}
\begin{pastaba}
Pagrindimas remiasi nagrinėjamų objektų apibrėžtimis ir tuo, kad kiekvienas veiksmas išlaiko pradinės situacijos prasmę. Formalus taikymas reikalauja nurodyti sąlygas ir patikrinti gautą rezultatą.
\end{pastaba}

\subsection{Taisyklės ir algoritmai}
\begin{enumerate}[label=\arabic*.]
\item Išskirkite duotus dydžius, sąvokas arba objektus.
\item Pasirinkite tinkamą užrašą: skaitinį reiškinį, formulę, lentelę, brėžinį, diagramą arba lygtį.
\item Atlikite veiksmus nuosekliai, kiekviename žingsnyje išlaikydami vienetus ir sąlygas.
\item Patikrinkite, ar gautas rezultatas atitinka uždavinio prasmę.
\end{enumerate}

\subsection{Iliustruojantys pavyzdžiai}
\begin{enumerate}[label=\arabic*.]
\item Pritaikykime temą paprastam skaičiavimui: pradinė reikšmė yra \(12\), pokytis yra \(3\).
\begin{align}
12+3&=15.
\end{align}
\textbf{Atsakymas: \(15\).}
\item Palyginkime du galimus rezultatus \(18\) ir \(24\).
\begin{align}
24-18&=6,\\
24&>18.
\end{align}
\textbf{Atsakymas: antras rezultatas didesnis \(6\) vienetais.}
\item Sudarykime ir išspręskime paprastą modelį, kai nežinomas dydis \(x\) su \(5\) sudaro \(20\).
\begin{align}
x+5&=20,\\
x&=15.
\end{align}
\textbf{Atsakymas: \(x=15\).}
\end{enumerate}

\subsection{Ryšys su kitomis temomis}
Ši tema papildo to paties koncentro skaičių, modelių, geometrijos arba duomenų temas ir padeda nuosekliai pereiti prie aukštesnių klasių uždavinių, kuriuose reikia derinti kelias matematikos sritis.


\section{Daugianariai ir faktorizavimas}\label{sec:daugianariai-ir-faktorizavimas}

\subsection{Apibrėžtys}
\begin{apibreztis}
Tema apima daugianarių veiksmus, greitosios daugybos formules, faktorizavimą ir pertvarkius. Ji nagrinėjama kaip atskira mokymo(si) turinio dalis pagal \emph{matematikaPrograma.pdf}.
\end{apibreztis}
\begin{apibreztis}
Matematinis modelis šioje temoje yra skaičių, simbolių, brėžinių, lentelių arba diagramų sistema, kuria aprašoma reali arba teorinė situacija.
\end{apibreztis}

\subsection{Teoremos ir savybės}
\begin{savybe}
Jeigu uždavinio sąlyga tiksliai perkeliama į matematinį modelį, tai modelio rezultatas gali būti interpretuojamas pradinėje situacijoje.
\end{savybe}
\begin{pastaba}
Pagrindimas remiasi nagrinėjamų objektų apibrėžtimis ir tuo, kad kiekvienas veiksmas išlaiko pradinės situacijos prasmę. Formalus taikymas reikalauja nurodyti sąlygas ir patikrinti gautą rezultatą.
\end{pastaba}

\subsection{Taisyklės ir algoritmai}
\begin{enumerate}[label=\arabic*.]
\item Išskirkite duotus dydžius, sąvokas arba objektus.
\item Pasirinkite tinkamą užrašą: skaitinį reiškinį, formulę, lentelę, brėžinį, diagramą arba lygtį.
\item Atlikite veiksmus nuosekliai, kiekviename žingsnyje išlaikydami vienetus ir sąlygas.
\item Patikrinkite, ar gautas rezultatas atitinka uždavinio prasmę.
\end{enumerate}

\subsection{Iliustruojantys pavyzdžiai}
\begin{enumerate}[label=\arabic*.]
\item Pritaikykime temą paprastam skaičiavimui: pradinė reikšmė yra \(12\), pokytis yra \(3\).
\begin{align}
12+3&=15.
\end{align}
\textbf{Atsakymas: \(15\).}
\item Palyginkime du galimus rezultatus \(18\) ir \(24\).
\begin{align}
24-18&=6,\\
24&>18.
\end{align}
\textbf{Atsakymas: antras rezultatas didesnis \(6\) vienetais.}
\item Sudarykime ir išspręskime paprastą modelį, kai nežinomas dydis \(x\) su \(5\) sudaro \(20\).
\begin{align}
x+5&=20,\\
x&=15.
\end{align}
\textbf{Atsakymas: \(x=15\).}
\end{enumerate}

\subsection{Ryšys su kitomis temomis}
Ši tema papildo to paties koncentro skaičių, modelių, geometrijos arba duomenų temas ir padeda nuosekliai pereiti prie aukštesnių klasių uždavinių, kuriuose reikia derinti kelias matematikos sritis.


\section{Lygčių sistemų modeliavimas}\label{sec:lygciu-sistemu-modeliavimas}

\subsection{Apibrėžtys}
\begin{apibreztis}
Tema apima realaus konteksto situacijų aprašymą lygčių sistemomis ir sprendinių interpretavimą. Ji nagrinėjama kaip atskira mokymo(si) turinio dalis pagal \emph{matematikaPrograma.pdf}.
\end{apibreztis}
\begin{apibreztis}
Matematinis modelis šioje temoje yra skaičių, simbolių, brėžinių, lentelių arba diagramų sistema, kuria aprašoma reali arba teorinė situacija.
\end{apibreztis}

\subsection{Teoremos ir savybės}
\begin{savybe}
Jeigu uždavinio sąlyga tiksliai perkeliama į matematinį modelį, tai modelio rezultatas gali būti interpretuojamas pradinėje situacijoje.
\end{savybe}
\begin{pastaba}
Pagrindimas remiasi nagrinėjamų objektų apibrėžtimis ir tuo, kad kiekvienas veiksmas išlaiko pradinės situacijos prasmę. Formalus taikymas reikalauja nurodyti sąlygas ir patikrinti gautą rezultatą.
\end{pastaba}

\subsection{Taisyklės ir algoritmai}
\begin{enumerate}[label=\arabic*.]
\item Išskirkite duotus dydžius, sąvokas arba objektus.
\item Pasirinkite tinkamą užrašą: skaitinį reiškinį, formulę, lentelę, brėžinį, diagramą arba lygtį.
\item Atlikite veiksmus nuosekliai, kiekviename žingsnyje išlaikydami vienetus ir sąlygas.
\item Patikrinkite, ar gautas rezultatas atitinka uždavinio prasmę.
\end{enumerate}

\subsection{Iliustruojantys pavyzdžiai}
\begin{enumerate}[label=\arabic*.]
\item Pritaikykime temą paprastam skaičiavimui: pradinė reikšmė yra \(12\), pokytis yra \(3\).
\begin{align}
12+3&=15.
\end{align}
\textbf{Atsakymas: \(15\).}
\item Palyginkime du galimus rezultatus \(18\) ir \(24\).
\begin{align}
24-18&=6,\\
24&>18.
\end{align}
\textbf{Atsakymas: antras rezultatas didesnis \(6\) vienetais.}
\item Sudarykime ir išspręskime paprastą modelį, kai nežinomas dydis \(x\) su \(5\) sudaro \(20\).
\begin{align}
x+5&=20,\\
x&=15.
\end{align}
\textbf{Atsakymas: \(x=15\).}
\end{enumerate}

\subsection{Ryšys su kitomis temomis}
Ši tema papildo to paties koncentro skaičių, modelių, geometrijos arba duomenų temas ir padeda nuosekliai pereiti prie aukštesnių klasių uždavinių, kuriuose reikia derinti kelias matematikos sritis.


\section{Tiesiniai ir netiesiniai sąryšiai I–II gimnazijos klasėse}\label{sec:tiesiniai-ir-netiesiniai-sarysiai-9-10}

\subsection{Apibrėžtys}
\begin{apibreztis}
Tema apima tiesinės, kvadratinės, atvirkštinio proporcingumo ir kitas funkcines priklausomybes. Ji nagrinėjama kaip atskira mokymo(si) turinio dalis pagal \emph{matematikaPrograma.pdf}.
\end{apibreztis}
\begin{apibreztis}
Matematinis modelis šioje temoje yra skaičių, simbolių, brėžinių, lentelių arba diagramų sistema, kuria aprašoma reali arba teorinė situacija.
\end{apibreztis}

\subsection{Teoremos ir savybės}
\begin{savybe}
Jeigu uždavinio sąlyga tiksliai perkeliama į matematinį modelį, tai modelio rezultatas gali būti interpretuojamas pradinėje situacijoje.
\end{savybe}
\begin{pastaba}
Pagrindimas remiasi nagrinėjamų objektų apibrėžtimis ir tuo, kad kiekvienas veiksmas išlaiko pradinės situacijos prasmę. Formalus taikymas reikalauja nurodyti sąlygas ir patikrinti gautą rezultatą.
\end{pastaba}

\subsection{Taisyklės ir algoritmai}
\begin{enumerate}[label=\arabic*.]
\item Išskirkite duotus dydžius, sąvokas arba objektus.
\item Pasirinkite tinkamą užrašą: skaitinį reiškinį, formulę, lentelę, brėžinį, diagramą arba lygtį.
\item Atlikite veiksmus nuosekliai, kiekviename žingsnyje išlaikydami vienetus ir sąlygas.
\item Patikrinkite, ar gautas rezultatas atitinka uždavinio prasmę.
\end{enumerate}

\subsection{Iliustruojantys pavyzdžiai}
\begin{enumerate}[label=\arabic*.]
\item Pritaikykime temą paprastam skaičiavimui: pradinė reikšmė yra \(12\), pokytis yra \(3\).
\begin{align}
12+3&=15.
\end{align}
\textbf{Atsakymas: \(15\).}
\item Palyginkime du galimus rezultatus \(18\) ir \(24\).
\begin{align}
24-18&=6,\\
24&>18.
\end{align}
\textbf{Atsakymas: antras rezultatas didesnis \(6\) vienetais.}
\item Sudarykime ir išspręskime paprastą modelį, kai nežinomas dydis \(x\) su \(5\) sudaro \(20\).
\begin{align}
x+5&=20,\\
x&=15.
\end{align}
\textbf{Atsakymas: \(x=15\).}
\end{enumerate}

\subsection{Ryšys su kitomis temomis}
Ši tema papildo to paties koncentro skaičių, modelių, geometrijos arba duomenų temas ir padeda nuosekliai pereiti prie aukštesnių klasių uždavinių, kuriuose reikia derinti kelias matematikos sritis.


\section{Trikampių teoremos ir įrodymai}\label{sec:trikampiu-teoremos-ir-irodymai}

\subsection{Apibrėžtys}
\begin{apibreztis}
Tema apima Pitagoro, Talio, sinusų ir kosinusų teoremų taikymą bei paprastus įrodymus. Ji nagrinėjama kaip atskira mokymo(si) turinio dalis pagal \emph{matematikaPrograma.pdf}.
\end{apibreztis}
\begin{apibreztis}
Matematinis modelis šioje temoje yra skaičių, simbolių, brėžinių, lentelių arba diagramų sistema, kuria aprašoma reali arba teorinė situacija.
\end{apibreztis}

\subsection{Teoremos ir savybės}
\begin{savybe}
Jeigu uždavinio sąlyga tiksliai perkeliama į matematinį modelį, tai modelio rezultatas gali būti interpretuojamas pradinėje situacijoje.
\end{savybe}
\begin{pastaba}
Pagrindimas remiasi nagrinėjamų objektų apibrėžtimis ir tuo, kad kiekvienas veiksmas išlaiko pradinės situacijos prasmę. Formalus taikymas reikalauja nurodyti sąlygas ir patikrinti gautą rezultatą.
\end{pastaba}

\subsection{Taisyklės ir algoritmai}
\begin{enumerate}[label=\arabic*.]
\item Išskirkite duotus dydžius, sąvokas arba objektus.
\item Pasirinkite tinkamą užrašą: skaitinį reiškinį, formulę, lentelę, brėžinį, diagramą arba lygtį.
\item Atlikite veiksmus nuosekliai, kiekviename žingsnyje išlaikydami vienetus ir sąlygas.
\item Patikrinkite, ar gautas rezultatas atitinka uždavinio prasmę.
\end{enumerate}

\subsection{Iliustruojantys pavyzdžiai}
\begin{enumerate}[label=\arabic*.]
\item Pritaikykime temą paprastam skaičiavimui: pradinė reikšmė yra \(12\), pokytis yra \(3\).
\begin{align}
12+3&=15.
\end{align}
\textbf{Atsakymas: \(15\).}
\item Palyginkime du galimus rezultatus \(18\) ir \(24\).
\begin{align}
24-18&=6,\\
24&>18.
\end{align}
\textbf{Atsakymas: antras rezultatas didesnis \(6\) vienetais.}
\item Sudarykime ir išspręskime paprastą modelį, kai nežinomas dydis \(x\) su \(5\) sudaro \(20\).
\begin{align}
x+5&=20,\\
x&=15.
\end{align}
\textbf{Atsakymas: \(x=15\).}
\end{enumerate}

\subsection{Ryšys su kitomis temomis}
Ši tema papildo to paties koncentro skaičių, modelių, geometrijos arba duomenų temas ir padeda nuosekliai pereiti prie aukštesnių klasių uždavinių, kuriuose reikia derinti kelias matematikos sritis.


\section{Statistinis patikimumas}\label{sec:statistinis-patikimumas}

\subsection{Apibrėžtys}
\begin{apibreztis}
Tema apima imties atsitiktinumą, dydį, duomenų patikimumą ir išvadų ribotumą. Ji nagrinėjama kaip atskira mokymo(si) turinio dalis pagal \emph{matematikaPrograma.pdf}.
\end{apibreztis}
\begin{apibreztis}
Matematinis modelis šioje temoje yra skaičių, simbolių, brėžinių, lentelių arba diagramų sistema, kuria aprašoma reali arba teorinė situacija.
\end{apibreztis}

\subsection{Teoremos ir savybės}
\begin{savybe}
Jeigu uždavinio sąlyga tiksliai perkeliama į matematinį modelį, tai modelio rezultatas gali būti interpretuojamas pradinėje situacijoje.
\end{savybe}
\begin{pastaba}
Pagrindimas remiasi nagrinėjamų objektų apibrėžtimis ir tuo, kad kiekvienas veiksmas išlaiko pradinės situacijos prasmę. Formalus taikymas reikalauja nurodyti sąlygas ir patikrinti gautą rezultatą.
\end{pastaba}

\subsection{Taisyklės ir algoritmai}
\begin{enumerate}[label=\arabic*.]
\item Išskirkite duotus dydžius, sąvokas arba objektus.
\item Pasirinkite tinkamą užrašą: skaitinį reiškinį, formulę, lentelę, brėžinį, diagramą arba lygtį.
\item Atlikite veiksmus nuosekliai, kiekviename žingsnyje išlaikydami vienetus ir sąlygas.
\item Patikrinkite, ar gautas rezultatas atitinka uždavinio prasmę.
\end{enumerate}

\subsection{Iliustruojantys pavyzdžiai}
\begin{enumerate}[label=\arabic*.]
\item Pritaikykime temą paprastam skaičiavimui: pradinė reikšmė yra \(12\), pokytis yra \(3\).
\begin{align}
12+3&=15.
\end{align}
\textbf{Atsakymas: \(15\).}
\item Palyginkime du galimus rezultatus \(18\) ir \(24\).
\begin{align}
24-18&=6,\\
24&>18.
\end{align}
\textbf{Atsakymas: antras rezultatas didesnis \(6\) vienetais.}
\item Sudarykime ir išspręskime paprastą modelį, kai nežinomas dydis \(x\) su \(5\) sudaro \(20\).
\begin{align}
x+5&=20,\\
x&=15.
\end{align}
\textbf{Atsakymas: \(x=15\).}
\end{enumerate}

\subsection{Ryšys su kitomis temomis}
Ši tema papildo to paties koncentro skaičių, modelių, geometrijos arba duomenų temas ir padeda nuosekliai pereiti prie aukštesnių klasių uždavinių, kuriuose reikia derinti kelias matematikos sritis.


\section{Tikimybiniai modeliai ir stebėjimas}\label{sec:tikimybiniai-modeliai-ir-stebejimas}

\subsection{Apibrėžtys}
\begin{apibreztis}
Tema apima klasikinius modelius, santykinį dažnį, bandymų kartojimą ir tikimybės įvertinimą stebėjimu. Ji nagrinėjama kaip atskira mokymo(si) turinio dalis pagal \emph{matematikaPrograma.pdf}.
\end{apibreztis}
\begin{apibreztis}
Matematinis modelis šioje temoje yra skaičių, simbolių, brėžinių, lentelių arba diagramų sistema, kuria aprašoma reali arba teorinė situacija.
\end{apibreztis}

\subsection{Teoremos ir savybės}
\begin{savybe}
Jeigu uždavinio sąlyga tiksliai perkeliama į matematinį modelį, tai modelio rezultatas gali būti interpretuojamas pradinėje situacijoje.
\end{savybe}
\begin{pastaba}
Pagrindimas remiasi nagrinėjamų objektų apibrėžtimis ir tuo, kad kiekvienas veiksmas išlaiko pradinės situacijos prasmę. Formalus taikymas reikalauja nurodyti sąlygas ir patikrinti gautą rezultatą.
\end{pastaba}

\subsection{Taisyklės ir algoritmai}
\begin{enumerate}[label=\arabic*.]
\item Išskirkite duotus dydžius, sąvokas arba objektus.
\item Pasirinkite tinkamą užrašą: skaitinį reiškinį, formulę, lentelę, brėžinį, diagramą arba lygtį.
\item Atlikite veiksmus nuosekliai, kiekviename žingsnyje išlaikydami vienetus ir sąlygas.
\item Patikrinkite, ar gautas rezultatas atitinka uždavinio prasmę.
\end{enumerate}

\subsection{Iliustruojantys pavyzdžiai}
\begin{enumerate}[label=\arabic*.]
\item Pritaikykime temą paprastam skaičiavimui: pradinė reikšmė yra \(12\), pokytis yra \(3\).
\begin{align}
12+3&=15.
\end{align}
\textbf{Atsakymas: \(15\).}
\item Palyginkime du galimus rezultatus \(18\) ir \(24\).
\begin{align}
24-18&=6,\\
24&>18.
\end{align}
\textbf{Atsakymas: antras rezultatas didesnis \(6\) vienetais.}
\item Sudarykime ir išspręskime paprastą modelį, kai nežinomas dydis \(x\) su \(5\) sudaro \(20\).
\begin{align}
x+5&=20,\\
x&=15.
\end{align}
\textbf{Atsakymas: \(x=15\).}
\end{enumerate}

\subsection{Ryšys su kitomis temomis}
Ši tema papildo to paties koncentro skaičių, modelių, geometrijos arba duomenų temas ir padeda nuosekliai pereiti prie aukštesnių klasių uždavinių, kuriuose reikia derinti kelias matematikos sritis.
